\chapter{Introdução}
\label{cap:introducao}

O avanço acelerado da Inteligência Artificial (IA) generativa tem impulsionado transformações profundas no cenário educacional, com os grandes modelos de linguagem (LLMs) emergindo como ferramentas promissoras para mediar a relação entre o estudante e o conhecimento \citeb{kasneci2023chatgpt}. Aplicações baseadas nesses modelos passaram a ser exploradas em tarefas como tutoria automatizada, sumarização de materiais didáticos e resposta a perguntas sobre conteúdos acadêmicos, abrindo novas possibilidades para o aprendizado assistido por IA. Contudo, apesar da alta capacidade de articulação textual, esses modelos apresentam uma limitação estrutural crítica para o ambiente educacional: a alucinação semântica. Esse fenômeno é caracterizado pela geração de respostas que, embora gramaticalmente fluidas, são factualmente incorretas ou completamente desconectadas das fontes de dados consultadas, o que compromete diretamente a confiabilidade da informação entregue ao estudante \citeb{raschka2024build}.

Para mitigar essa vulnerabilidade e ancorar as respostas em fontes de dados verificáveis, \citeonline{NEURIPS2020_6b493230} introduziram a arquitetura de \textit{Retrieval-Augmented Generation} (RAG), uma abordagem que atua como um intermediário que extrai blocos de informação de uma base de conhecimento externa e os injeta no contexto do modelo antes da geração da resposta, garantindo a proveniência dos dados e reduzindo a ocorrência de alucinações. No contexto educacional, essa arquitetura representa um avanço relevante, pois possibilita que sistemas de IA respondam com base no conteúdo real dos materiais didáticos fornecidos pelo próprio estudante ou pela instituição.

Entretanto, o RAG tradicional foi concebido para operar sobre repositórios textuais lineares, o que representa uma limitação significativa diante da natureza dos documentos acadêmicos reais. Conforme apontado por \citeonline{gao2023retrieval}, o processamento textual padrão frequentemente fragmenta e corrompe dados semiestruturados presentes em PDFs. Além disso, mecanismos de recuperação baseados puramente em texto mostram-se insuficientes para lidar com documentos cuja transmissão de conhecimento depende de layouts complexos, tabelas, infográficos e diagramas visuais, elementos amplamente presentes em livros didáticos e artigos científicos \citebdois{faysse2025colpali}{yun2025lilac}. Nesse contexto, a evolução para arquiteturas de \textit{Multimodal Retrieval-Augmented Generation} (MRAG) impõe-se como caminho natural, uma vez que a compreensão integral de um documento educacional exige que o motor de busca seja capaz de interpretar simultaneamente as dimensões textuais e visuais que estruturam a informação.

No entanto, apesar dos avanços que a transição para MRAG representa, as soluções presentes na literatura ainda divergem entre eficiência computacional e precisão granular, não atendendo plenamente às exigências do domínio educacional. Abordagens baseadas puramente em Visão Computacional, como a arquitetura ColPali, proposta por \citeonline{faysse2025colpali}, introduziram a recuperação direta de imagens de páginas inteiras utilizando mecanismos de \textit{``interação tardia''}. Apesar de serem mais rápidas por eliminarem processos custosos de reconhecimento óptico de caracteres (\textit{Optical Character Recognition} -- OCR), essas arquiteturas operam em uma granularidade fixa de página, o que frequentemente insere ruído visual irrelevante no contexto do modelo e dilui a precisão em consultas específicas \citeb{yun2025lilac}.

Para contornar esse excesso de ruído e mitigar alucinações, frameworks recentes como o SLEUTH \citeb{liu2025resolving}, empregam múltiplos agentes de IA colaborativos para refinar a busca em documentos longos. Paralelamente, sistemas como o LILaC \citeb{yun2025lilac} fragmentam o documento em grafos de componentes em camadas para interligar explicitamente parágrafos, tabelas e imagens, alcançando alta precisão em raciocínios que exigem cruzar informações de múltiplas páginas. Contudo, tanto a orquestração de múltiplos agentes quanto a construção de grafos complexos combinada com a decomposição de consultas via LLMs impõem uma elevada carga computacional e aumentam a latência do sistema \citebdois{liu2025resolving}{yun2025lilac}.

Além disso, o foco primário dessas tecnologias tem sido o domínio aberto, científico ou administrativo \citebdois{liu2025resolving}{yun2025lilac}, negligenciando a estrutura pedagógica única, expressa na relação direta entre teoria, exercícios e diagramas explicativos, intrínseca aos materiais didáticos.

O presente trabalho, intitulado EDGAR, propõe o desenvolvimento de uma arquitetura de Geração Aumentada por Recuperação Multimodal (MRAG), concebida para o contexto educacional. O sistema tem como objetivo indexar e recuperar, com alta precisão e baixo custo computacional, tanto o conteúdo textual quanto a estrutura visual de documentos didáticos complexos, incluindo tabelas, gráficos e diagramas. Em vez de depender de uma estrutura custosa baseada em múltiplos agentes ou da modelagem de grafos genéricos de domínio aberto, a solução introduz um método de indexação voltado à compreensão da hierarquia semântica do material de estudo. Ao isolar e recuperar o recorte visual exato associado ao seu respectivo contexto textual, a arquitetura busca fornecer ao modelo gerador um bloco de informação mais consistente, contextualizado e com menor presença de ruídos.

Dessa maneira, este projeto apresenta contribuições relevantes para a promoção da segurança, da veracidade e da confiabilidade da informação no processo de ensino-aprendizagem mediado por IA, uma vez que a ocorrência de alucinações compromete diretamente o propósito educacional da tecnologia. Nesse sentido, a relevância do trabalho está relacionada à preservação do elo pedagógico entre elementos textuais e visuais, mantendo-os alinhados à estrutura original do conteúdo. Para possibilitar o uso, a demonstração e a avaliação prática da arquitetura proposta, este trabalho prevê o desenvolvimento do LuminIA Reader, uma aplicação cliente baseada na web concebida para permitir a submissão de documentos, o envio de perguntas e a visualização das respostas e evidências retornadas pela arquitetura. Como resultado, a abordagem busca mitigar falhas de forma proativa por meio da curadoria estrutural do contexto, reduzindo ruídos visuais e textuais antes da etapa de geração da resposta pelo modelo de linguagem.

Para alcançar os objetivos propostos, o desenvolvimento da arquitetura EDGAR seguirá um percurso metodológico organizado em etapas de especificação arquitetural, modelagem dos \textit{pipelines} de dados, definição dos componentes tecnológicos e planejamento da avaliação experimental. Os recursos previstos baseiam-se no uso de ferramentas e modelos de código aberto (\textit{open-source}), incluindo redes neurais pré-treinadas para análise de layout, mecanismos de extração textual, modelos de descrição visual e bancos de dados vetoriais voltados à indexação e recuperação semântica. Essa escolha busca favorecer a reprodutibilidade da pesquisa e a viabilidade de implementação em ambientes computacionais acessíveis, com apoio de infraestrutura remota para as etapas de maior custo computacional.

A avaliação da arquitetura proposta será conduzida por meio de um plano experimental destinado a analisar tanto o desempenho computacional dos fluxos de processamento quanto a qualidade semântica das respostas geradas. Serão consideradas métricas relacionadas à latência no fluxo síncrono, ao tempo de processamento no fluxo assíncrono, à precisão na recuperação de evidências, à relevância das respostas e à fidelidade ao contexto recuperado. Como limitação deste estágio da pesquisa, o escopo restringe-se à análise da arquitetura EDGAR em documentos de domínio educacional normalizados no formato PDF. Embora a proposta possa futuramente ser adaptada para outros formatos documentais, esta etapa concentra-se em materiais estruturados em páginas, como apostilas, capítulos de livros e apresentações convertidas para PDF. Portanto, não são contemplados, neste momento, formatos de texto fluido sem paginação nativa, como arquivos EPUB brutos, nem documentos baseados em hipertextos interativos externos.

\section{Objetivos}
\label{sec:objetivos}

\subsection{Objetivo Geral}
\label{subsec:objetivo_geral}

Desenvolver uma arquitetura de Geração Aumentada por Recuperação Multimodal (MRAG) voltada à consulta de materiais de estudo em PDF, combinando segmentação de layout e extração de elementos visuais estruturantes, com o objetivo de fornecer ao modelo gerador evidências mais específicas e contextualizadas, contribuindo para a redução de alucinações semânticas no contexto educacional.

\subsection{Objetivos Específicos}
\label{subsec:objs_especificos}

\begin{itemize}
    \item Projetar \textit{pipelines} de processamento e indexação multimodal para identificar elementos textuais e visuais, como tabelas, gráficos e diagramas, vinculando-os de forma granular ao contexto documental correspondente;

    \item Desenvolver mecanismos de refinamento contextual para selecionar os componentes mais relevantes e reduzir ruídos visuais e textuais do material didático antes da etapa de geração da resposta;

    \item Desenvolver um protótipo da aplicação cliente LuminIA Reader, concebido como meio de interação com a arquitetura EDGAR, permitindo a submissão de documentos, o envio de perguntas e a visualização das respostas e evidências recuperadas;

    \item Avaliar a viabilidade técnica da arquitetura por meio de testes experimentais, utilizando métricas automatizadas e cenários baseados em materiais didáticos reais para analisar a qualidade da recuperação, a fidelidade das respostas e a latência computacional dos fluxos propostos.
\end{itemize}